
This work establishes three synthetic proof-of-concept contributions pending real data validation:

\begin{enumerate}
\item \textbf{Integration framework}: Hierarchical Bayesian calibration integrates consistency-based and conformal prediction methods through mutual calibration, achieving expected calibration error of 0.043 (below 0.05 threshold) while maintaining 92\% coverage in synthetic proof-of-concept validation. Prior work applies these methods independently or in simple cascades. Bayesian joint calibration improves both signals beyond independent application in synthetic experiments.
\end{enumerate}

\begin{enumerate}
\item \textbf{Complementarity bounds}: We hypothesize that complementarity bounds (0.3 < $\rho$ < 0.7) determine when joint calibration adds value. Synthetic validation with controlled correlation (target $\rho$=0.5, achieved $\rho \approx$ 0.43--0.46) demonstrates the mechanism. Observed correlation occupies this range in synthetic validation. If $\rho$ > 0.8, methods are redundant and single-method optimization is preferred. If $\rho$ < 0.2, methods are independent and cascade is optimal. At moderate $\rho$, joint calibration may be optimal.
\end{enumerate}

\begin{enumerate}
\item \textbf{Computational efficiency mechanism}: Hierarchical Bayesian calibration reduces cost by 30\% (2,800 vs. 4,000 forward passes per 1,000 queries) compared to conformal-only through consistency-informed weighting in synthetic proof-of-concept validation. Epistemic structure may guide statistical calibration to achieve efficiency without sacrificing coverage.
\end{enumerate}

The correlation stability finding in synthetic validation supports the complementarity hypothesis. Correlation remains $\rho \approx$ 0.43 across synthetic tasks where correlation was controlled (epistemic-heavy, aleatoric-heavy, mixed), suggesting the mechanism operates as designed in proof-of-concept validation.

\subsubsection{Implications and Future Directions}

The complementarity framework (quantify correlation $\rho$, determine integration strategy) applies beyond consistency and conformal methods. Any pair of uncertainty estimators can be analyzed for complementarity.

For practitioners, hierarchical Bayesian calibration may enable deployment in production systems requiring both statistical rigor and computational efficiency, pending real data validation to confirm synthetic results.

The method may resolve the paradox between impossibility results (absolute hallucination detection requires oracle) and empirical success of consistency methods. Consistency methods measure epistemic structure, not absolute truth, while conformal methods provide aleatoric bounds. These are complementary dimensions, not competing truth claims.

Three research directions follow:

\begin{enumerate}
\item \textbf{Real data validation}: Conduct full-scale validation on TruthfulQA (817 samples), HH-RLHF (8,552 samples), SQuAD (11,873 samples) to confirm correlation stability ($\rho \approx$ 0.43) and performance metrics (expected calibration error = 0.043) beyond synthetic validation.
\end{enumerate}

\begin{enumerate}
\item \textbf{Pure epistemic/aleatoric tasks}: Test on closed-book question answering (pure epistemic uncertainty) and subjective classification (pure aleatoric uncertainty) to validate whether correlation remains $\rho \approx$ 0.43. If correlation persists, robust complementarity is confirmed. If correlation varies, task-dependent correlation provides design guidance.
\end{enumerate}

\begin{enumerate}
\item \textbf{Few-shot domain adaptation}: Explore whether consistency priors enable conformal calibration with fewer than 100 labeled examples. The 30\% cost reduction in synthetic validation suggests calibration set size requirements may be relaxed.
\end{enumerate}

\begin{enumerate}
\item \textbf{Multi-modal extension}: Apply hierarchical Bayesian calibration to vision-language models where epistemic uncertainty (visual grounding failures) and aleatoric uncertainty (image ambiguity) may exhibit similar complementarity.
\end{enumerate}

\subsubsection{Honest Limitations (Revisited)}

We reiterate three principled limitations:

\begin{enumerate}
\item \textbf{Synthetic proof-of-concept}: Validation used synthetic data (200 samples per dataset) with controlled correlation (target $\rho$=0.5). Real model inference on full-scale datasets required for production deployment. The core mechanism is validated; specific metrics require real data confirmation.
\end{enumerate}

\begin{enumerate}
\item \textbf{Labeled calibration requirement}: Hierarchical Bayesian calibration requires labeled examples. Few-shot adaptation and zero-shot deployment remain open challenges.
\end{enumerate}

\begin{enumerate}
\item \textbf{Out-of-distribution untested}: Domain shift experiments not conducted. Out-of-distribution detection claims remain speculative. Core calibration contribution stands independently.
\end{enumerate}

These limitations define the scope of the validated contribution while pointing to natural next steps.

\subsubsection{Closing}

Uncertainty quantification methods for large language models need not operate in isolation. By recognizing potential complementarity between consistency-based and statistical approaches---quantifying their moderate correlation in synthetic validation and designing hierarchical Bayesian integration to exploit it---we provide a synthetic proof-of-concept demonstrating simultaneous improvement in calibration quality and computational efficiency. This resolves the tradeoff between statistical rigor and practical deployment in synthetic experiments, pending real data validation to confirm performance at scale.
